% ID: SEL-D0039-OB-E03
% ESTADO: cotejada. Segunda pasada visual y borrador compilado revisados.
% Fuente: Selectividad/matematicas_aplicadas/Ua4MATEMÁTICAS APLICADAS examen.pdf
% SHA-256: 78ff6bdd48cbac71e6cc47d6ddf3dcb9492540ecb9cbe2032ebfec02e39fb66d
% Página: 2; opción B; ejercicio original 3.

La proporción de personas de una población que tiene una determinada enfermedad es de 1 por cada 500 personas. Se dispone de una prueba para detectar dicha enfermedad. La prueba detecta la enfermedad en el 90\% de los casos en que la persona está enferma, pero también da como enfermas al 5\% de las personas sanas.
\begin{enumerate}[label=\alph*)]
\item \textbf{(1.25 puntos)} Se elige al azar una persona y se le hace la prueba. \textquestiondown{}Cuál es la probabilidad de que haya sido diagnosticada correctamente?
\item \textbf{(1.25 puntos)} Si la prueba ha diagnosticado que la persona está enferma, \textquestiondown{}cuál es la probabilidad de que realmente lo esté? \textquestiondown{}Y de que esté sana?
\end{enumerate}
